\section*{Chapter \ref{ch:qm:testing-and-multiple-testing} --- Testing and Multiple Testing}

\begin{solution}{exo:qm:testing-and-multiple-testing:1}
$p = 1 - \Phi(2.5) = 0.0062$; Bonferroni over twenty variants, $20 \times 0.0062 = 0.124$: not significant at 5\%.
\end{solution}

\begin{solution}{exo:qm:testing-and-multiple-testing:2}
Size $1 - \Phi(1.645) = 5\%$; power $\P(\mathcal N(2, 1) > 1.645) = \Phi(0.355) = 0.639$.
\end{solution}

\begin{solution}{exo:qm:testing-and-multiple-testing:3}
Bonferroni's threshold is 0.0125: it rejects the first two. Holm compares the ordered $p$-values with 0.0125, 0.0167, 0.025 and
0.05 in turn: 0.001, 0.012 and 0.02 pass and 0.3 does not, so it makes three rejections. The adjusted $p$-values are, for
Bonferroni, 0.004, 0.048, 0.08 and 1; for Holm, 0.004, 0.036, 0.04 and 0.3.
\end{solution}

\begin{solution}{exo:qm:testing-and-multiple-testing:4}
$T = ((1.645 + 1.282)/0.7)^2 = 17.5$ years.
\end{solution}

\begin{solution}{exo:qm:testing-and-multiple-testing:5}
With known zero mean, $\mathrm{LR} = n\ln\hat s^2 - n_1\ln\hat s_1^2 - n_2\ln\hat s_2^2$ with the pooled $\hat s^2 = (1 + 1.44)/2 = 1.22$ (in
$\%^2$): $1\,000\ln1.22 - 500\ln1.44 = 16.53$. Against $\chi^2_1$, $p = 4.8 \times 10^{-5}$: the variance changed.
\end{solution}

\begin{solution}{exo:qm:testing-and-multiple-testing:6}
$1 - \Phi(2)^{10} = 1 - 0.97725^{10} = 0.206$. The formula gives $(1 - \gamma_E)\Phi^{-1}(0.9) + \gamma_E\Phi^{-1}(1 - 1/(10e)) = 1.57$
(the exact expected maximum is 1.54).
\end{solution}

\begin{solution}{exo:qm:testing-and-multiple-testing:7}
With 500 screens (\texttt{fdr\_correlated()}): the \omterm{def:qm:testing-and-multiple-testing:fdr}{false discovery rate} is 3.0\%, within the 5\% bound (positive dependence), with
about 63 discoveries per screen. But the realised proportion is volatile: it exceeds 10\% in 7.8\% of screens, when the common
factor pushes all the null statistics up together. The procedure controls an average, not each screen.
\end{solution}

\begin{solution}{exo:qm:testing-and-multiple-testing:8}
Twelve tests at 5\% produce 0.6 false positives on average; three or more occur with probability 2.0\% if the signal is useless
everywhere and the markets independent, so the count is mild evidence of something. The flaw is the selection: the three
winners were chosen because their backtests were good, so their in-sample performance is biased upward (the winner's curse),
and the markets are not independent. Report the family result, and estimate performance on data not used to choose the markets.
\end{solution}

\begin{solution}{pb:qm:testing-and-multiple-testing:1}
\textbf{1.} Lookback 42 days: $t = 3.11$, annualised \omterm{def:qm:estimation:sharpe}{Sharpe ratio} 0.98.
\textbf{2.} $p = 1 - \Phi(3.11) = 0.00094$.
\textbf{3.} 33 of 200.
\textbf{4.} 1.04: every rule trades the same history, so its trending episodes lift the whole family together.
\textbf{5.} 0.95 with the 43-day rule; 0.28 with the 5-day and 0.20 with the 204-day rule.
\textbf{6.} $200 \times 0.00094 = 0.187$; critical value $\Phi^{-1}(1 - 0.05/200) = 3.48$.
\textbf{7.} 0.187: Holm multiplies the smallest $p$-value by $m$, as Bonferroni does; it gains only further down the list.
\textbf{8.} 0.076: the step-up takes the minimum of $p_{(k)}m/k$ over the ranks above, and the 33 small $p$-values make those small.
\textbf{9.} $1 - (1 - 0.00094)^{200} = 0.171$.
\textbf{10.} 0.030 from 200\,000 Gaussian draws with the estimated correlation; the 5\% critical value of the maximum is 2.92.
\textbf{11.} $\ln(1 - 0.030)/\ln(1 - 0.00094) = 32$ trials.
\textbf{12.} Three.
\textbf{13.} In 2.6\% of them, close to the simulated 3.0\% (the sampling error of 2\,000 histories is 0.4 points).
\textbf{14.} The probability that the true \omterm{def:qm:estimation:sharpe}{Sharpe ratio} is positive is 0.999. Deflated for 200 trials: benchmark 0.87 a year, \omterm{def:qm:testing-and-multiple-testing:dsr}{deflated
Sharpe ratio} 0.63; for 32 trials: 0.67 and 0.84.
\textbf{15.} Barely: with Student-$t_4$ returns the best rule reaches 3.11 in 2.35\% of 2\,000 fresh histories; the $t$-statistics are
still close to normal after 2\,520 days.
\textbf{16.} A 5\% test rejects a true null 5\% of the time, and the best rule's adjusted $p$-value is 3\%; besides, the history is the
336th drawn, a search the correction was not told about.
\textbf{17.} A holdout never used in the search, an economic reason for trend at that horizon, performance net of costs, and an
adjusted $p$-value that counts every family the team tried.
\textbf{18.} The number of variants and families tried, their correlation, the adjusted $p$-values and the method, the sample period and
what was fixed before looking at the data.
\textbf{19.} Named result: \emph{the best of two hundred}: the best of two hundred correlated momentum rules, naive $p = 0.00094$, has a
max-statistic adjusted $p$-value of 0.030 (Bonferroni: 0.187); the search amounts to 32 independent trials.
\textbf{20.} It protects against the searches it is told about, and against nothing the researcher did not count.
\end{solution}

\begin{solution}{iq:qm:testing-and-multiple-testing:1}
The probability, if the null is true, of a statistic at least as extreme as the one observed. It is not the probability that the null
is true, not the probability that the result will repeat, and not a measure of the size of the effect.
\iqlookfor{The conditional direction ($\P(\text{data} \mid H_0)$, not $\P(H_0 \mid \text{data})$), and that effect size needs its own number.}
\end{solution}

\begin{solution}{iq:qm:testing-and-multiple-testing:2}
Probably not. If all were noise and independent, the best of 100 exceeds 2.8 with probability $1 - \Phi(2.8)^{100} = 0.23$;
correlation lowers that, but then ask how many variants and choices came before the hundred. Adjust for the search (max-statistic or
Bonferroni), then test on a holdout.
\iqlookfor{The one-in-four arithmetic, the effect of correlation, and a holdout as the real test.}
\end{solution}

\begin{solution}{iq:qm:testing-and-multiple-testing:3}
Bonferroni or Holm (family-wise error) when one false positive is costly, such as the one strategy that goes live; Benjamini--Hochberg
(\omterm{def:qm:testing-and-multiple-testing:fdr}{false discovery rate}) when screening many candidates for a combined model, where a few false ones among many true ones are acceptable.
Under strong dependence use the max-statistic version or Benjamini--Yekutieli.
\iqlookfor{Matching the error rate to the cost of a false discovery; Holm dominates Bonferroni.}
\end{solution}

\begin{solution}{iq:qm:testing-and-multiple-testing:4}
The annual estimate has \omterm{def:qm:estimation:estimator}{standard error} about $1/\sqrt T$, so $t = \sqrt T$: four years for $t = 2$, and at 5\% size with 80\% power
$(1.645 + 0.842)^2 = 6.2$ years. Autocorrelated returns and fat tails make it longer; a search makes it longer still.
\iqlookfor{The $1/\sqrt T$ rule and the power calculation, not only $t = 2$.}
\end{solution}

\begin{solution}{iq:qm:testing-and-multiple-testing:5}
Use the law of the maximum rather than the union bound: estimate the correlation of the variants' P\&L and simulate the maximum of
correlated normals, or bootstrap the maximum jointly (\omterm{def:qm:testing-and-multiple-testing:rc}{Reality Check}, SPA, Romano--Wolf step-down). Report the effective number of
trials. Bonferroni is valid but very conservative here.
\iqlookfor{Max-statistic over union bound, joint resampling that keeps the dependence, and counting families as well as variants.}
\end{solution}

\begin{solution}{iq:qm:testing-and-multiple-testing:6}
The implicit multiplicity of choices that depend on the data even when one analysis is run. In machine learning: feature choices made after
looking at validation scores, early stopping and hyperparameters tuned on the test set, reruns with new seeds until the curve looks good,
and a test set reused across many model iterations until it becomes a training set.
\iqlookfor{Concrete leak points, and the discipline of a holdout touched once.}
\end{solution}

\begin{solution}{iq:qm:testing-and-multiple-testing:7}
Little. The test is weak in the tails, which is where return distributions differ; with parameters fitted to the same data its
tabulated critical values are too lenient (Lilliefors), so it rejects even less often; and failing to reject is not acceptance. Compare
tail quantiles or exceedance counts directly, with simulated critical values.
\iqlookfor{Weak tail power, the fitted-parameter correction, and absence of evidence versus evidence of fit.}
\end{solution}
